\documentclass{article}
\usepackage[paperwidth=595.3bp,paperheight=841.9bp,margin=0bp]{geometry}
\usepackage{fontspec,eso-pic,graphicx,amsmath,booktabs,array}
\pagestyle{empty}
\setlength{\parindent}{0pt}
\setlength{\unitlength}{1bp}
\newcommand{\FontPath}{/Users/jinana/Library/Containers/com.kingsoft.wpsoffice.mac/Data/.kingsoft/office6/data/fonts/online_ext_font_cache/}
\defaultfontfeatures{Ligatures=NoCommon}
\setmainfont{Times New Roman}
\newfontfamily{\Times}{Times New Roman}
\newfontfamily{\Italic}{Times New Roman Italic}
\newfontfamily{\Bold}{Times New Roman Bold}
\newfontfamily{\Song}{HYShuSongErKW.ttf}[Path=\FontPath]
\newfontfamily{\SongHeavy}{HYShuSongErKW.ttf}[Path=\FontPath,FakeBold=2.86]
\newfontfamily{\Hei}{HYZhongHeiKW.ttf}[Path=\FontPath,FakeBold=2.86]
\newfontfamily{\Kai}{HYKaiTiKW.ttf}[Path=\FontPath]
\newfontfamily{\KaiHeavy}{HYKaiTiKW.ttf}[Path=\FontPath,FakeBold=2.86]
\newfontfamily{\Deng}{HYZhongDengXianKW.ttf}[Path=\FontPath]
\newcommand{\ReportPage}[1]{%
  \AddToShipoutPictureFG*{\AtPageLowerLeft{\setlength{\unitlength}{1bp}#1}}%
  \null\newpage
}
\newcommand{\Block}[4]{%
  \put(#1,\dimexpr841.9bp-#2bp\relax){\makebox(0,0)[lt]{%
  \begin{minipage}[t]{#3bp}\fontsize{10bp}{16bp}\selectfont
  \setlength{\parskip}{5bp}#4\end{minipage}}}%
}

\usepackage{xeCJK}
\usepackage{fix-cm}
\setCJKmainfont{HYShuSongErKW.ttf}[Path=\FontPath]
\begin{document}
\ReportPage{%
  \Block{90}{101}{415.3}{%
    1. 电压换算\par
    \input{parameters.tex}
    \[
      H=\frac{N_1}{LR_1}U_X,\qquad
      B=\frac{R_2C}{N_2S}U_Y .
    \]
    \vspace{-7bp}
    \begin{center}
    \fontsize{9bp}{13bp}\selectfont
    \setlength{\tabcolsep}{4bp}\input{conversion.tex}\par
    表中换算式的 \(U_X,U_Y\) 均以 mV 为单位。
    \end{center}
  }%
  \Block{90}{257}{415.3}{%
    2. 饱和磁滞回线（上）与基本磁化曲线（下）\par
    \vspace{-5bp}
    \includegraphics[width=415.3bp]{curves.pdf}
  }%
  \Block{90}{502}{415.3}{%
    3. 剩磁和矫顽力\par
    \vspace{-3bp}
    在过轴相邻点间作线性插值，求 \(H=0\) 时的 \(B_{r\pm}\) 和
    \(B=0\) 时的 \(H_{c\pm}\)；两支绝对值取平均，得 \(B_r\)、\(H_c\)。\par
    \vspace{-3bp}
    {\centering\fontsize{8.5bp}{12bp}\selectfont
    \setlength{\tabcolsep}{4bp}\input{results.tex}\par}
  }%
  \Block{90}{631}{415.3}{%
    两样品均呈闭合磁滞回线，\(H=0\) 时仍有剩磁，需反向磁场才能使
    \(B=0\)；基本磁化曲线随 \(H\) 增大逐渐趋缓。样品2的剩磁和矫顽力
    均较大，回线更宽；样品1更易退磁。\par
    主要误差来自示波器光标读数及零点偏移、过轴点间的插值近似，
    以及 \(R_1,R_2,C,L,S\) 的参数偏差。有限的 \(R_2C\) 会使积分近似
    带来幅值和相位误差；退磁不充分也会使回线产生偏移。
  }%
}
\ReportPage{%
  \Block{111}{247}{394.3}{%
    主要根据矫顽力及回线形状判断：软磁材料的矫顽力小、回线窄，
    易磁化也易退磁；硬磁材料的矫顽力大、回线宽，磁化后不易退磁，
    适于制作永磁体。\par
    按指导书 \(H_c\leq1000\,\mathrm{A/m}\) 的软磁判据，
    本次两种样品均表现为软磁材料。
  }%
  \Block{111}{507}{394.3}{%
    铁磁材料有磁滞，撤去外磁场后仍可能保留剩磁，使下一次测量受
    原有磁化历史影响。退磁时施加方向交替、幅值逐渐减小的磁场，
    使磁滞回线逐步收缩，最终使宏观剩磁趋近于零。\par
    具体操作：保持交流励磁，先顺时针缓慢调大信号源幅度，使样品
    达到磁饱和；再逆时针缓慢、连续调小幅度，直到示波器上的回线
    缩为一点、励磁幅度降至零。不能用突然断电代替逐渐减小励磁。
  }%
}
\end{document}
